\sect{Natural Language Facilitated Human-Robot Cooperation}{language}

NL Facilitated Human-Robot Cooperation (NLC) refers to the use of natural language as the primary medium for communication and coordination between humans and robots working toward shared objectives. Unlike traditional control interfaces that rely on structured commands, programming, or gesture vocabularies, NLC enables intuitive interaction by allowing users to express intentions through unrestricted linguistic input. This approach integrates human cognitive reasoning and goal formulation with the robot’s physical execution capabilities, forming a more natural interface for collaborative and adaptive systems \cite{doi:10.1177/1729881419851402}. 

Earlier NLC frameworks typically followed multi-stage symbolic pipelines for interpreting and executing user instructions. Human utterances were first parsed into formal semantic structures and then translated into executable actions using predefined ontologies or rule-based mappings \cite{LAURIA2002171}. These architectures demonstrated the feasibility of language-guided control but were constrained by limited vocabulary coverage, rigid grammar, and the absence of contextual understanding. The symbolic nature of these models required explicit task definitions and was prone to failure when faced with linguistic ambiguity or unseen phrasing \cite{doi:10.1177/1729881419851402}.

Advances in data-driven learning have substantially improved language understanding within robotic cooperation frameworks. Probabilistic and logic-based models introduced partial generalization by associating linguistic features with task parameters and sensor-derived context \cite{doi:10.1177/1729881419851402}. However, these methods still required handcrafted mappings and domain-specific training, restricting scalability and flexibility in real-world applications. As robotic systems grow more complex, the need for higher-level abstraction and open-ended intent interpretation has motivated the integration of language models capable of reasoning beyond fixed symbolic structures.

Recent developments in LLMs address many of these limitations by enabling semantic comprehension, contextual reasoning, and intent recognition from unstructured input \cite{wang2025function}. Unlike grammar-based or probabilistic systems, LLMs generalize across linguistic variability, capturing task dependencies and causal relations within a single inference process. This capability eliminates much of the manual engineering traditionally required for language grounding and allows robots to interpret, plan, and execute tasks directly from NL goals. Building on these advancements, LLMs with structured function-calling capabilities now bridge NL intent and executable robotic actions, marking a shift from symbolic interpretation to unified agentic reasoning.

